% =============================================================================
% theme.tex — workspace-wide Beamer theme selection.
%
% Every deck under ./slides/ loads THIS file (via \InputIfFileExists). Change it
% once here and all decks in this workspace restyle together. If you delete this
% file, decks fall back to the bundled Simple theme.
%
% --- Option 1: use the bundled Simple theme (default) ---------------------------
% To tweak its colors, set them BEFORE \usetheme (all use \providecolor):
%   \definecolor{simpleAccent}{RGB}{160,40,40}   % primary accent
%   \definecolor{simpleAlert}{RGB}{158,108,24}   % caveat / alert blocks
%   \definecolor{simpleExample}{RGB}{35,132,104} % example / desired blocks
% The Simple theme also accepts options:
%   \usetheme[eyebrow={PAPER READING DECK}]{Simple}  % small title-page label
%   \usetheme[density=dense]{Simple}                 % tighter margins + type
%   \usetheme[divider=off]{Simple}                   % no section divider frames
%   \usetheme[overflowguard=on]{Simple}              % error if a slide overflows
% Full guide: beamerthemeSimple-manual.md. Demo deck: demo-theme.tex.
\usetheme{Simple}

% --- Option 2: use any Beamer theme installed in your TeX distribution ----------
% Comment out the line above and uncomment one of these:
%   \usetheme{metropolis}
%   \usetheme{Madrid}
%   \usetheme{Berlin}

% --- Option 3: bring your own custom theme --------------------------------------
% Drop your `beamerthemeMyTheme.sty` (and any images it needs) at this workspace
% root — the deck Makefile's TEXINPUTS=../../ makes it discoverable — then use:
%   \usetheme{MyTheme}
% If your theme references image assets, put them at the workspace root or in the
% deck's figures/ (both are on \graphicspath); or set an absolute path.
% =============================================================================
